\chapter{Weber-Kräfte und das Quantenpotential Q}

Die Weber-Kraft ist instantan und geschwindigkeitsabhängig. Ihre allgemeine Form lautet:

\begin{equation}
\vec{F} = \frac{Q_1 Q_2}{4\pi \epsilon_0 r^2} \left\{ \left[ 1 - \frac{\dot{r}^2}{c^2} + \beta_{\text{Kraft}} \frac{r\ddot{r}}{c^2} \right] \hat{r} + \frac{2(\dot{r} \cdot \vec{v})}{c^2} \vec{v} \right\}. \tag{5.1}
\end{equation}

Es gibt drei Ausprägungen:

\begin{itemize}
    \item \textbf{WED} (Weber-Elektrodynamik): \( \beta = 2 \)
    \item \textbf{WG} (Weber-Gravitation für Massen): \( \beta = 0,5 \)
    \item \textbf{WG für Photonen}: \( \beta = 1 \) (kompensiert durch Q)
\end{itemize}

\section{Die Weber-Kräfte im Detail}

\subsection{Die Weber-Elektrodynamik (WED)}

Für die Elektrodynamik gilt \( \beta_{\text{WED}} = 2 \):

\begin{equation}
\vec{F}_{\text{WED}} = \frac{q_1 q_2}{4\pi \epsilon_0 r^2} \left\{ \left[ 1 - \frac{\dot{r}^2}{c^2} + 2\frac{r\ddot{r}}{c^2} \right] \hat{r} + \frac{2(\dot{r} \cdot \vec{v})}{c^2} \vec{v} \right\}. \tag{5.2}
\end{equation}

Die WED erfüllt \emph{actio = reactio} und erhält den Impuls. Sie ist reversibel.

\subsection{Die Weber-Gravitation (WG) für Massen}

Für massive Körper (Planeten, Sterne) gilt \( \beta_{\text{WG}} = 0,5 \):

\begin{equation}
\vec{F}_{\text{WG}} = -\frac{GMm}{r^2} \left\{ \left[ 1 - \frac{\dot{r}^2}{c^2} + \frac{1}{2}\frac{r\ddot{r}}{c^2} \right] \hat{r} + \frac{2(\dot{r} \cdot \vec{v})}{c^2} \vec{v} \right\}. \tag{5.3}
\end{equation}

Die WG ist irreversibel – sie erzeugt die DSTT-Drift (\( \dot{\beta} > 0 \)).

\subsection{Die Weber-Gravitation für Photonen}

Für Photonen (masselose Teilchen) gilt \( \beta = 1 \):

\begin{equation}
\vec{F}_{\text{WG}}^{(\gamma)} = -\frac{GMm}{r^2} \left\{ \left[ 1 - \frac{\dot{r}^2}{c^2} + \frac{r\ddot{r}}{c^2} \right] \hat{r} + \frac{2(\dot{r} \cdot \vec{v})}{c^2} \vec{v} \right\}. \tag{5.4}
\end{equation}

Hier greift das Quantenpotential Q ein. Es kompensiert den \( \beta = 1 \)-Wert so, dass effektiv \( \beta = 0,5 \) übrig bleibt.

\section{Die Instantaneität der Weber-Kräfte und von Q}

Die Weber-Kräfte (WG und WED) sind \emph{instantan}. Sie wirken direkt zwischen den Teilchen – sie sind eine Fernwirkung im Newtonschen Sinne, aber mit geschwindigkeitsabhängiger Struktur (IIR-Rekursion). 

Das Quantenpotential Q ist ebenfalls \emph{instantan}. Es wirkt nicht-lokal auf das gesamte System, ohne dass ein Signal mit \(c\) unterwegs sein muss.

\subsection{Der Unterschied zwischen Weber-Kräften und Q}

Der Unterschied liegt \emph{nicht} in der Instantaneität, sondern in der \emph{Reichweite} und \emph{Struktur}:

\begin{itemize}
    \item \textbf{WG/WED:} Lokale, direkte Zwei-Körper-Wechselwirkung. Sie wirken nur zwischen den beteiligten Teilchen.
    \item \textbf{Q:} Globale, nicht-lokale Führungsgröße, die von der gesamten Dichteverteilung \( \rho \) abhängt.
\end{itemize}

Beide sind instantane Aspekte derselben Informationswelle der IWT.

\section{Das Quantenpotential Q als instantane Welle}

Das Quantenpotential in der de Broglie-Bohm-Theorie lautet:

\begin{equation}
Q = -\frac{\hbar^2}{2m} \frac{\nabla^2 \sqrt{\rho}}{\sqrt{\rho}}. \tag{5.5}
\end{equation}

Q ist \emph{nicht-lokal} und \emph{instantan}. Es ist die globale Führungsgröße, die das gesamte System beeinflusst.

\subsection{Die Wellennatur von Q}

Q ist keine Materie, kein Teilchen, kein Feld im klassischen Sinne. Es ist eine \emph{Krümmung der Wahrscheinlichkeitsdichte} – ein Muster im Informationsraum. 

Als instantane Welle hat Q folgende Eigenschaften:

\begin{itemize}
    \item Es ist \emph{überall gleichzeitig} – es braucht keine Zeit, um sich auszubreiten.
    \item Es ist \emph{nicht-lokal} – es wirkt auf alle Teilchen gleichzeitig.
    \item Es ist \emph{kein Signal} – es ist ein Muster im Informationsnetzwerk.
\end{itemize}

\subsection{Die Verbindung zur fraktalen Energieverteilung}

Die fraktale Energieverteilung (Kapitel 2) ist der \emph{zeitliche Mittelwert} der Q-Welle. Die "Löcher" sind die Minima der Q-Welle, die "Berge" sind die Maxima.

\section{Q als Regler des Kreislaufs}

Q erfüllt drei Funktionen:

\begin{enumerate}
    \item \textbf{Materieentstehung:} \( \sigma_{\text{cre}} \propto |\nabla Q|^2 \).
    \item \textbf{Verteilung der Energie:} instantane, nicht-lokale Umverteilung.
    \item \textbf{Rückkopplung:} Aufnahme der Rotverschiebungsenergie.
\end{enumerate}

\subsection{Die Identitätsbeziehung}

Die vollständige WDBT-Kraft ist:

\begin{equation}
\vec{F}_{\text{WDBT}} = \vec{F}_{\text{WG}} + \vec{F}_Q, \qquad \vec{F}_Q = -\nabla Q. \tag{5.6}
\end{equation}

Es gilt die exakte Identitätsbeziehung:

\begin{equation}
\frac{|\vec{F}_{\text{WDBT}}^{\text{tan}}|}{|\vec{F}_{\text{WDBT}}^{\text{rad}}| + |\vec{F}_{\text{WDBT}}^{\text{tan}}|} \equiv \beta_{\text{DSTT}}(t). \tag{5.7}
\end{equation}

Die Weber-Kräfte und Q sind zwei Seiten derselben instantanen Welle – der Informationswelle der IWT.

\section{Zusammenfassung}

\begin{enumerate}
    \item Die Weber-Kraft ist instantan und geschwindigkeitsabhängig.
    \item Es gibt drei Ausprägungen: WED (\( \beta = 2 \)), WG für Massen (\( \beta = 0,5 \)), WG für Photonen (\( \beta = 1 \), kompensiert durch Q).
    \item Q ist das Quantenpotential – eine instantane, nicht-lokale Welle.
    \item WG/WED und Q sind \emph{beide} instantan – der Unterschied liegt in der Reichweite (lokal vs. global) und der Struktur.
    \item Q ist der Regler des Kreislaufs: es erzeugt Materie, verteilt Energie und nimmt Rotverschiebungsenergie auf.
    \item Die Weber-Kräfte und Q sind zwei Seiten derselben instantanen Welle – der Informationswelle der IWT.
    \item Die Identitätsbeziehung (5.7) verbindet die Weber-Kräfte mit der DSTT-Drift.
\end{enumerate}